\subsection*{同号异内容如何留下可处置事实}

EV1 的 3000 分输入已经提交，后来同 EV1 携带 3500 分到达。仅捕获主键异常再 ACK，会把不同内容混成普通重投。这里保留原 I1、F1 与原结果，把来件完整输入、原输入、原因和负责角色写入核查项。再次收到同一冲突，恢复同一项，避免每次投递再开一张无人处理的单。

核查完成以前，业务尚未接受这份变动内容。核查者查明发布错误，可以保存决定并要求来源以新的、有依据的业务事实更正，或者在已有授权决定的条件下执行有依据的修复。已提交的 EV1 输入保持可追查，后续处理通过决定与修复记录接上原事实。

同一效果也要比较内容。EV2 与 EV1 是不同事件，若两者都报告 O1、P1/C1/X1、Q7/CNY 3000 和相同商品数量，就可各自保存接收结果，引用同一 F1。第九章新 R 的事实唯一约束会阻止正常路径为同一 C1 再建一个事件。这里的双事件历史用于检查旧路径已有其他编号、或上游重复创建事件时，履约效果受哪项约束保护。若 EV2 报告同一 O1 却换了商品或金额，效果唯一冲突触发核查。仅比较 B 就恢复成功，会把这份异参输入藏进原履约单。

\CTenFigure{ch10-conflict-route}{接收输入冲突与效果输入冲突分别比较原绑定，并把未决责任共同保存}{fig:ch10-conflict-route}

退款、补发和重新履约具有新的业务原因。退款 RF1 沿退款授权与资金身份处理。经批准的补发 AR1 可以形成 \code{(O1,reship,AR1)}，并验证原履约事实、补发数量和审批依据。它们需要各自的处理规则和唯一关系。本章 \code{initial} 的受控入口只接受既定初次授权。补发入口取得 AR1 的批准依据以后，再按自己的身份保存效果与结果。

\subsection*{升级专题：已经成功的 O1 怎样补建通知}

前两节已经让新系统完成接收身份 I 与效果身份 B 的共同仲裁。现在把源端与消费端一起放回升级现场：第八章已有成功 O1，旧 R 没有 Outbox；旧服务还可能发送过另一种格式的通知。上线第九章的新 R，只会改变后来执行的事务，已经结束的成功也要有明确的迁移位置。

为这批旧成功记录安排一次受控补建。切换时先停止旧 R 入口，等待在途本地事务结束，再把当时已经成功的 C 及其原事实摘要登记为持久的旧记录集合。这个集合也称为 cohort，即本次切换明确圈定的一批记录。登记由迁移程序保存，补建者据它判断哪些成功事实属于旧规则。全体补建完成并核对以后，才开放只走新规则的 R。原结果的只读查询可以继续，尚未决的 C 后来取得成功时直接走新 R，平台付款身份保持原值。暂停、登记与补建屏障共同划出可核查的边界，每项旧成功都有明确的补建归属。

\CTenFigure{ch09-migration-barrier}{先圈定旧成功集合并完成补建，再开放新 R，形成可恢复的源端切换屏障}{fig:ch09-migration-barrier}

补建事务按 C1→T1 锁定原关联，读取旧集合中的登记，确认 C1 已成功、P1 已 \code{paid}、T1 已 \code{done}，并核对唯一 O1 与保存结果中的 P/C/O/X、币种和金额一致。然后从原购买输入、报价快照和已保存支付事实构造通知，按同一事实唯一约束查找事件。尚无事件才插入，已有事件则逐字段比较并恢复。事件与本项补建完成共同提交。并发补建者在原 C 处串行，退出以后按登记继续。零行、冲突或原事实摘要不符，均整笔回滚并登记核查。补建只增加发布依据，原付款、订单、结果和任务终态保持原值。

若旧服务曾直接发送过通知，补建后的消息可能再次报告 O1 已付款。新分配的事件编号只能统一新路径的副本。旧副本若没有这个编号，下游就要沿 O1 对应的履约动作认出相同效果。启用这批补发以前，旧、新通知应进入同一项履约效果约束。对于发送情况不明的旧通知，迁移记录保存未知结论和负责核查的归属。

新规则启用后，受控补建只面向已登记的旧成功集合。已经按新 R 提交的成功记录若缺少事件，普通结果重放报告完整性异常，并保留核查归属。核查者据原提交及发送证据判断是否误删、是否已有消息到达，再决定恢复办法。普通重放径直分配一个新事件，会把这些情况混在一起。


\subsection*{消费端屏障：旧通知与补建事件怎样汇合}

源端 cohort 已经说明哪些旧成功需要补建 EV。消费端另有自己的屏障。先停旧消费者，排空本地事务，并处理已经发出的外部动作。旧执行可能晚到时，使其继续受同一业务身份保护，或取得能关闭旧执行资格的证据。若只把新程序上线，旧程序继续按原入口发货，新程序查到本地尚无 F 后又建立一份工作，两个入口就可能分别执行同一初次履约。

对每个登记的 O 保存原事实摘要与迁移依据。已经存在旧 F 的，核对 O、原商品数量、金额、P/C/X 等效果输入，导入原 F 及其原始来源，同时导入已有回报依据，或按同一 F 补齐唯一回报待办。迁移登记保留旧系统编号、原始记录、核对结论与负责角色，导入的 \code{origin=\allowbreak legacy-import} 和 \code{origin\_evidence} 指向这份依据。导入事务按 B 保存授权与原 F，再共同保存回报关联。以后 EV1 到达，消费者就能把其效果输入与这张已导入 F 逐字段比较，并保存真实发生的本次接收结果。

旧程序可能只有外部发货记录，履约库缺少 F，也可能发送过请求而结果不明。后者把 B 的授权行置为 \code{hold}，建立负责核查的持久项，阻止自动建立新履约单。核查从仓储、承运商及原调用身份取得证据：已经发货则恢复原外部映射，导入原 F、已知的 \code{dispatched} 进度与回报，调度者据此跳过初次发货。只证明旧建单完成时，则保留其已知建单阶段。确认尚未执行且旧在途资格已关闭，或所有迟到请求受同一身份保护，才允许沿原 B 继续。若批准真正的新补发，使用 AR1 的新动作并保留原未知项的处置关系。

旧通知如果没有 \code{event\_id}，可在明确的旧规则下使用 \code{(legacy-order-paid,O1)} 作为消息种类与消息号：该旧格式每份订单只表达一次已付款初次履约，原完整快照能够从来件或受信迁移登记取得并比较。格式不满足这个条件时，入口保存原件进入核查。能够恢复其原业务含义以后，再按有依据的身份处理。

EV1 的新接收身份与旧通知不同，两者经同一 B 恢复原 F。反过来，先收到旧通知、它在新受控入口首次建 F，后来 EV1 到达，也只保存第二份接收结果。源端补建屏障和消费端迁移屏障完成各自责任以后，补发才有这条共同效果依据。

\CTenFigure{ch10-original04}{源端补建与消费端迁移是两道屏障，未知旧执行阻止自动新履约}{fig:ch10-migration}

\subsection*{乱序需要哪些额外信息}

对于同一事件的重复，I 的相等足以找到原接收记录。再加入不同事实，问题会变化。设查询投影先收到 O1 的已发货快照，稍后收到较早的待调度快照。两个事件号都没有处理过，去重表允许两次修改，后一次盲目覆盖就会使界面从已发货退回待调度。

即使单个赋值重复两次不改变结果，交错历史也要检查。先写已发货、再写待调度，结果与业务先后相反~\cite{handling-duplicate-messages}。状态名也未必构成一个能简单取最大值的序列。取消、退款和补发各有自己的条件：退款取决于原资金与退款授权，补发取决于原履约和新的批准依据。消费者按相应状态迁移规则判断，再保存结论与其依据。

本章当前主方案接收第九章的 EV1，它没有对象业务序号。其允许动作是为已付款初次授权建立 F，迟到的同授权事件只恢复 F，已经 \code{dispatched} 的 F 保留当前进度。因此这个特定动作可以在允许乱序的通知路径上处理。若输入改成多步增量迁移，就要使每步迁移能找到自己的前驱。

一种额外方案是在源端对象事务内分配 \code{object\_sequence}，使同一 O 的连续业务变更取得 1、2、3。这个序号标识对象业务变更的顺序，与发布尝试号及代理投递编号分开。消费者保存连续完成位置 \(h\)，以及尚未应用的完整来件。处理普通增量时，先锁定本地同一对象的状态与 \(h\)，核对前驱，再把状态变化、连续位置和等待项的处理决定共同提交。\(h=1\) 时先收到 3，把 3 持久保存为等待项，继续位置为 1。取得并完成 2 后，才能依规则执行 3，再共同推进到 3。

等待项若已完整持久保存并有可扫描的后续工作，就可以确认原投递，确认的含义是责任已转交，业务效果尚未完成。单纯把 3 的编号写进最大已见位置并丢弃 2，会使增量永久缺失。保存最大编号的简化需要额外的对象语义，例如较晚的完整快照确实覆盖较早状态~\cite{handling-duplicate-messages}。

2 迟迟未到，消费者沿原对象和序号补取源日志，或者由负责者核查产生与发布依据。源端本来就没有产生 2、日志已清理以及 2 暂时迟到，分别对应不同恢复动作。对可以从源事实重新计算的查询投影，还可以补取快照。源端在同一对象状态边界保存快照及覆盖版本 \(k\)，说明这份状态已经完整包含 1 至 \(k\) 的业务变更。若导出状态与版本分别读取，版本 5 却配到变更 4 的旧状态，消费者会误以为第 5 项已经包含，后续也就失去了补上它的依据。

快照本身也可能迟到。设补取请求发出后，普通增量已把本地推进到 \(h=6\)，较早导出的 \(k=5\) 快照才返回。直接安装会丢掉第 6 项内容；只保留数字 6 而覆盖状态，则让连续位置与实际内容不符。消费者先核验快照来源、对象及覆盖版本，再进入与普通增量相同的对象行锁，读取锁定后的当前 \(h\)。

当 \(k<h\) 时，保存旧快照的处置观察，保留当前内容与位置。当 \(k=h\) 时，核对同一版本的投影内容，一致就返回已保存的同版本状态；投影缺失或损坏时，可以根据已经核验的同版本快照恢复内容，并保存重建依据，含义冲突则交给核查。当 \(k>h\) 时，确认快照覆盖整个前缀，才安装内容并把 \(h\) 推到 \(k\)。例如本地 \(h=1\)，可信快照完整覆盖至 5，就可安装后从 6 继续。

安装事务同时决定哪些等待项已被快照覆盖，以及哪些大于 \(k\) 的增量要继续处理。投影状态、\(h\)、等待项状态及可扫描的后续工作共同提交。普通增量取得同一对象锁后重读 \(h\)，避免快照与增量相互覆盖。若采用条件版本更新代替行锁，安装和增量都校验自己读取的同一对象版本，冲突时回滚重读，再决定上述分支。这样，安装者先读到 \(h=1\)、增量随后推进到 6 的历史，也会在仲裁处重新比较。这里恢复的是可重建投影，付款、发货等外部动作继续沿原业务身份核对效果和结果，快照安装只修改投影数据。

同一份 \(k=5\) 的可信快照，在锁内读到不同 \(h\) 时会有不同动作。图~\ref{fig:ch10-snapshot-admission} 分别取 \(h=6\)、5、1，展示旧快照处置、同版本核对与向前安装的区别。

\CTenFigure{ch10-snapshot-admission}{快照安装与普通增量共用对象仲裁，锁内比较覆盖版本并共同保存后续责任}{fig:ch10-snapshot-admission}

\CTenFigure{ch10-original05}{重复、状态倒退与额外有序方案中的缺口具有不同恢复依据}{fig:ch10-ordering}

\subsection*{重放和升级沿哪个范围继续}

查询投影因解析错误需要重建，可以建立新范围，例如 \code{order.query.v2:rebuild-1}，从原事件和有据快照重新计算查询表。新投影拥有新的输出位置，读取与重建都作用于查询表。切换查询入口时，再根据覆盖版本和增量追赶状态决定何时使用它。

履约处理器从程序构建 17 升到 18，继续使用 \code{fulfillment.initial.v1}。重新播放 EV1 时，先恢复原 I1。即使误换了接收范围，B1 的效果唯一也保护同一初次履约。删除 F 和授权记录再换消费组，会删除这项保护，因此回收与历史保留需要第十二章的专门规则。

事件 \code{schema\_version} 决定怎样读原输入，程序版本决定运行哪份代码。新程序保留原版本读取与比较路径。暂时不能解释某个旧版本时，持久保存原件并登记兼容处理任务，待相应读取器恢复后继续处理。若直接使用默认新版本解析，原输入的含义就可能发生变化。

修复确实需要追加业务效果时，以有授权的修复操作保存决定、对象范围和结果。原 F、原事件及其处理结果成为修复依据，追加动作从自己的授权入口执行。恢复快照也需要保留原身份和效果证据。把数据恢复到缺少 I/F 的过去，却允许旧副本重新进入，已经改变了本章已提交事实持续保留的前提，第三部分继续处理这类历史。
